\documentclass[10pt,a4paper]{amsart}
\usepackage[margin=19mm]{geometry}
\usepackage{amsmath,amssymb,mathtools}
\usepackage[scr=rsfs,cal=euler]{mathalfa}
\usepackage{xfrac}
\usepackage[unicode,hidelinks,bookmarks=false]{hyperref}
\newcommand{\ie}{i.e.\ }
\newcommand{\cf}{cf.\ }
\newcommand{\resp}{respectively\ }
\newcommand{\Subsection}{\subsection}
\providecommand{\nobreakdash}{\nobreak\hbox{-}}
\setlength{\emergencystretch}{2em}
\multlinegap0pt
\usepackage{bm}
\usepackage{bbm}
\usepackage{dsfont}
\usepackage{xspace}
\usepackage{cancel}
\usepackage{mathtools}
\usepackage{amsthm}
\makeatletter
\@ifundefined{thm}{\newtheorem{thm}{Thm}}{}
\@ifundefined{lem}{\newtheorem{lem}[thm]{Lemma}}{}
\makeatother
\newcommand{\Pic}{{\rm Pic}}
\newcommand{\StructureSheaf}[1]{{\mathcal O}_{#1}}
\newcommand{\hide}[1]{}
\title[Cycles on abelian 2n-folds of Weil type from secant sheaves ]{Cycles on abelian 2n-folds of Weil type from secant sheaves on abelian n-folds}
\author{Eyal Markman}
\date{Source: arXiv:2502.03415v2 (CC BY 4.0)}
\begin{document}
\maketitle

\noindent\emph{Source and licence.} Cycles on abelian 2n-folds of Weil type from secant sheaves on abelian n-folds. Version arXiv:2502.03415v2 (CC BY 4.0), distributed under the Creative Commons Attribution 4.0 International licence, as recorded in the arXiv licence metadata for this exact version and shown on the abstract page https://arxiv.org/abs/2502.03415v2. Full rights notice: licenses/arxiv-2502.03415-CC-BY-4.0.txt.\par\smallskip

\noindent\textit{Editorial scope.} Counted unit: the Lem whose source label is \texttt{lemma-intersection-of-C-and-Sigma-has-length-at-most-2} in the article (source0.tex lines 5475--5501, source label \texttt{lemma-intersection-of-C-and-Sigma-has-length-at-most-2}), delivered together with the article's own complete original proof of it (source0.tex lines 5503--5733, one \texttt{proof} environment of 231 lines). No mathematics is written, repaired or re-derived: statement and proof are the article's own text line for line. Required context retained uncounted: lemma-two-translates-of-C (lines 5740--5743); lemma-intersection-is-either-empty-or-of-length-2 (lines 5759--5765). Every definition, display and label of those lines is retained unchanged. Omitted from this package and not redistributed: 1--5474, 5502--5502, 5734--5739, 5744--5758, 5766--10494. Nothing inside a retained excerpt is omitted or altered: every retained body is delivered exactly as the article's lines are written, so no omission receipt of any kind accompanies this package. Citations: the retained excerpts carry no citation command at all, so no bibliography entry of the article is transcribed and no reference is redistributed. Reference adaptation: none was needed and none was made; every reference inside the delivered unit resolves inside it, so no unresolved reference or citation is produced. Printed slips kept literal: no printed word, symbol, hypothesis or number of the counted statement or of its proof was altered. Layout and class: the article's own class is \texttt{amsart} with packages including longtable, float, amssymb, amsbsy, amscd, amsmath, amsthm, bbold, mathrsfs, url, graphicx; the wrapper is the lane's pinned \texttt{amsart} editorial wrapper at 10pt on A4, which restates the theorem-like environments the retained text needs and adds the article's own macro definitions that the retained text needs (\texttt{\textbackslash Pic}, \texttt{\textbackslash StructureSheaf}, \texttt{\textbackslash hide}), with extra wrapper packages (bm, bbm, dsfont, xspace, cancel, mathtools). The delivered numbering is the unit's own. Assets: no retained span contains a figure environment, a graphics-inclusion command, a PDF-page inclusion, a table or a TikZ picture; no image file is copied into this package, the package declares no asset dependency and no image file is redistributed with it. Licence: version arXiv:2502.03415v2 of this article is distributed under the Creative Commons Attribution 4.0 International licence, as recorded in the arXiv licence metadata for this exact version and shown on the abstract page https://arxiv.org/abs/2502.03415v2. A full rights notice is delivered at licenses/arxiv-2502.03415-CC-BY-4.0.txt. Characters: every retained excerpt is pure ASCII, so the delivered engine, XeLaTeX with its default Latin Modern text font, reports no missing character.
\begin{lem}
\label{lemma-two-translates-of-C}
Two distinct translates of $AJ(C)$ in $\Pic^1(C)$ intersect along a subscheme of length at most $1$.
\end{lem}

\begin{lem}
\label{lemma-intersection-is-either-empty-or-of-length-2}
Let $t$ be a point of $\Pic^0(C)$. If $C_p\cap \tau_t(\Sigma_p)$ is non-empty, then 
$2p-t\sim a+b$, for a unique effective divisor $a+b$, $a,b\in C$. In that case 
the union $Z_{p,t}:=C_p\cup\tau_t(\Sigma_p)$ is contained in each of 
$\tau_{p-a}(\Theta)$ and $\tau_{p-b}(\Theta)$. 
\end{lem}

\begin{lem}
\label{lemma-intersection-of-C-and-Sigma-has-length-at-most-2}
\begin{enumerate}
\item
\label{lemma-item-length-at-most-2}
Given $t\in \Pic^0(X)$, the intersection $C_i\cap \tau_t(\Sigma_j)$ is either empty, or a subscheme of length $2$.
\item
\label{lemma-item-length-is-2-iff}
The intersection subscheme $C_i\cap \tau_t(\Sigma_j)$ has length $2$, if and only if 
there exists $u\in\Pic^0(C)$, such that $C_i\cup \tau_t(\Sigma_j)$ is contained in $\tau_u(\Theta)$. 
\item
\label{lemma-item-one-or-two-theta-translates}
If  $C_i\cup \tau_t(\Sigma_j)$ is contained in $\tau_u(\Theta)$, then 
one of the following holds.
\begin{enumerate}
\item
\label{case-a-of-intersection-lemma}
The union $C_i\cup \tau_t(\Sigma_j)$ is the complete intersection $\tau_u(\Theta)\cap\tau_{u'}(\Theta)$ 
of a unique pair of translates of $\Theta$. 
Furthermore, $(\tau_u\iota\tau_{-u})(C_i)\neq\tau_t(\Sigma_j)$.
\item
\label{case-b-of-intersection-lemma}
The divisor $\tau_u(\Theta)$ is the unique translate of $\Theta$ containing $C_i\cup \tau_t(\Sigma_j)$. 
Furthermore, $(\tau_u\iota\tau_{-u})(C_i)=\tau_t(\Sigma_j)$ and the canonical line bundle of $C_i\cup \tau_t(\Sigma_j)$ is the restriction of $\StructureSheaf{X}(2\tau_u(\Theta))$.
\end{enumerate}
\end{enumerate}
\end{lem}

\begin{proof}
Given $p\in C$, set $C_p:=C_{AJ(p)}$.  We may assume, without loss of generality, that $C_i=C_p$, for some $p\in C$. 
Let $s$ be a point of $\Pic^0(C)$, such that $\tau_t(\Sigma_j)=\tau_s(\Sigma_p)$.
%Let $\iota:X\rightarrow X$ be the involution given by $\iota(L)=\omega_C\otimes L^{-1}$. 


\underline{Step 0:} We prove part (\ref{lemma-item-one-or-two-theta-translates})  with the exception of the uniqueness in part (\ref{case-a-of-intersection-lemma}) which is postponed to Step 4 below.
Observe  the equivalence 
\begin{equation}
\label{equivalence-of-inclusions}
C_p\cup \tau_s(\Sigma_p)\subset \tau_u(\Theta)
\Leftrightarrow
C_p\cup \tau_s(\Sigma_p)\subset \tau_{s-u}(\Theta).
\end{equation}
Indeed, the right inclusion is obtained by applying $\tau_s\circ\iota$ to both sides of the left inclusion and using the equalities $\iota(\Theta)=\Theta$ and $\tau_s\circ\iota=\iota\circ\tau_{-s}$.

If $s\neq 2u$, then the union $C_p\cup \tau_s(\Sigma_p)$ is contained in the two distinct translates $\tau_u(\Theta)$ and
$\tau_{s-u}(\Theta)$ and so the union must be their complete intersection, as the cohomology classes of both are equal. If $\tau_u(\Theta)$ is the unique translate of $\Theta$ containing $C_p\cup \tau_s(\Sigma_p)$, then
%\footnote{We will see in Step 4 that in fact $s=0$ and $u=0$.} 
$s=2u$ and so $(\tau_u\iota\tau_{-u})(C_p)=\tau_{2u}(\iota(C_p))=\tau_s(\Sigma_p)$. 
If there exist two distinct translates of $\Theta$ containing $C_p\cup \tau_s(\Sigma_p)$, then the latter is their complete intersection, by the equality of the cohomology classes. 

Assume that $(\tau_u\iota\tau_{-u})(C_p)=\tau_s(\Sigma_p)$. Then $\tau_{2u}(\Sigma_p)=\tau_s(\Sigma_p)$. Hence, $s=2u$.
We will see 
%in Step 4(b) that this implies $s=u=0$ and 
in Step 3 that $\tau_u(\Theta)$ is the unique translate of $\Theta$ containing both $C_p$ and $\tau_{2u}(\Sigma_p)$. This would prove the inequality in part \ref{case-a-of-intersection-lemma}.

In both cases $Z':=C_i\cup\tau_t(\Sigma_j)$ is a divisor on a smooth surface $\tau_u(\Theta)$, and so its canonical sheaf is a line bundle.
In case (\ref{case-a-of-intersection-lemma}) the canonical line bundle $\omega_{Z'}$  is the restriction of
$\StructureSheaf{X}(\tau_u(\Theta)+\tau_{u'}(\Theta))$, as the normal bundle is the restriction of 
$\StructureSheaf{X}(\tau_u(\Theta))\oplus \StructureSheaf{X}(\tau_{u'}(\Theta))$.
Hence, $H^0(Z',\omega_{Z'}(-\tau_u(\Theta)-\tau_{u'}(\Theta)))$ is one dimensional. Case (\ref{case-b-of-intersection-lemma})
is a limit case and we conclude that $H^0(Z',\omega_{Z'}(-2\tau_u(\Theta)))$ does not vanish, by semi-continuity. The line bundle 
$\omega_{Z'}(-2\tau_u(\Theta))$ restricts to each irreducible component of $Z'$ as a line bundle of degree $0$, hence the existence of a non-zero global section implies that it is the trivial line-bundle.

\underline{Step 1:}
Let $q\in C$ be a point not equal to $p$ and set $y=p-q$. We prove first the equalities
\begin{eqnarray}
\label{eq-intersection-is-union-of-two-curves}
\Theta\cap \tau_y(\Theta)&=&C_p\cup \Sigma_q,
\\
\label{eq-C-p-and-Sigma-q-intersect-at-two-points}
C_p\cap\Sigma_q&=&\{\StructureSheaf{C}(p+r),\StructureSheaf{C}(p+t)\},
\end{eqnarray}
such that $\omega_C\cong\StructureSheaf{C}(p+q+r+t)$. The points $r$ and $t$ are the two other points on the intersection of the line through $\varphi_{\omega_C}(p)$ and $\varphi_{\omega_C}(q)$ with the canonical curve $\varphi_{\omega_C}(C)$. The curves $C_p$ and $\Sigma_q$ are tangent at $\StructureSheaf{C}(p+r)$, if $r=t$. 

Proof of (\ref{eq-intersection-is-union-of-two-curves}):
It suffices to prove the inclusion $(C_p\cup \Sigma_q)\subset (\Theta\cap \tau_y(\Theta))$, as the cohomology classes of both sides are equal. 
The curves $C_p$ and $C_q$ are contained in $\Theta$ and  $\iota(\Theta)=\Theta$, hence $\Sigma_p$ and $\Sigma_q$ are  contained in $\Theta$.
Now $C_q$ is contained in $\Theta$ and so $C_p=\tau_y(C_q)$ is contained in $\tau_y(\Theta)$. 
%Hence, $C_p$ is contained in $\Theta\cap \tau_y(\Theta)$.
In addition, $\iota(\tau_{-y}(\Theta))=\tau_y(\iota(\Theta))=\tau_y(\Theta)$, and so 
\begin{equation}
\label{eq-Sigma-q-is-translate-of-Sigma-p-by-p-q}
\Sigma_q=\iota(C_q)=\iota(\tau_{-y}(C_p))=\tau_y(\iota(C_p))=\tau_y(\Sigma_p)
\end{equation}
is contained in  $\tau_y(\Theta).$

%Observe next that $C_p\cap\Sigma_q$ consists of the two points $\StructureSheaf{C}(p+r)$ and $\StructureSheaf{C}(p+t)$,
%Indeed, 
Proof of (\ref{eq-C-p-and-Sigma-q-intersect-at-two-points}):
\[
p+r\in\Sigma_q\Leftrightarrow \exists t\in C, \ \omega_C(-q-t)\cong\StructureSheaf{C}(p+r)
\Leftrightarrow 
\exists t\in C, \  \omega_C\cong\StructureSheaf{C}(p+q+r+t).
\]
If $r=t$ the curves $C_p$ and $\Sigma_q$ are tangent at $\StructureSheaf{C}(p+r)$, since their intersection number in $\Theta$ is $2$.

\underline{Step 2:}
We show next that given $y\in\Pic^0(C)$, the curve $C_p$ is contained in $\tau_y(\Theta)$, if and only if $y=\StructureSheaf{C}(p-r)$, for some $r\in C$. 
\begin{equation}
\label{eq-criterion-for-inclusion}
C_p\subset \tau_y(\Theta) \Leftrightarrow \exists r\in C, \ y=\StructureSheaf{C}(p-r).
\end{equation}
The ``if'' implication is clear. We prove the ``only if'' implication. Assume that $C_p$ is contained in $\tau_y(\Theta)$.
Let $L$ be the line bundle represented by $-y$. Then $\tau_{-y}(C_p)$ is contained in $\Theta$, if and only if the line bundle $\StructureSheaf{C}(p+q)\otimes L$ is effective, for all $q\in C$.  In this case the point $p$ will belong to
the support of the effective divisor in the linear system $|\StructureSheaf{C}(p+q)\otimes L|$, for some $q$ as we vary $q$ in $C$.
Hence, if $C_p$ is contained in $\tau_y(\Theta)$, then $L\cong \StructureSheaf{C}(r-q)$, for some points $q, r\in C$.
Furthermore, $\StructureSheaf{C}(r-q+p+q')$ is effective, for all $q'\in C$. So $q$ is a base point of
$\StructureSheaf{C}(r+p+q')$, for the generic $q'\in C$ for which $h^0(\StructureSheaf{C}(r-q+p+q'))=1$ (i.e., for $q'\in C$, such that the images of $r, p, q'$ on the canonical curve are not colinear, or if $p=r$, then for $q'\in C$, such that $q'$ is not on the line tangent to the canonical curve at $p$). 
For such generic $q'$ the base locus is $\{r, p, q'\}$. 
Hence $q$ is equal to one of $r$ or $p$. If $q=p$ we are done. If $q=r$, then $L$ is trivial and we are done. 

Applying $\iota$ to (\ref{eq-criterion-for-inclusion}) and replacing $y$ by $-y$ we get
\begin{equation}
\label{eq-Sigma-p-is-contained-in-theta-translate-iff}
\Sigma_p\subset \tau_{y}(\Theta) \Leftrightarrow \exists r\in C, \ y=r-p.
\end{equation}

Combing (\ref{eq-criterion-for-inclusion}) and (\ref{eq-Sigma-p-is-contained-in-theta-translate-iff}) we see that $\Theta$ is the only translate of $\Theta$ containing both $C_p$ and $\Sigma_p$. Indeed, if $p-r=r'-p$, for points $r,r'\in C$, then $2p=r+r'$ and so $p=r=r'$, since $C$ is non-hyperelliptic.

\underline{Step 3:}
Combining the results of the above two steps we conclude that if $\Theta\cap\tau_t(\Theta)=C_p\cup C'$, for some $t\in \Pic^0(C)$ and some curve $C'$, then $t=p-q$, for some $q\in C$, and $C'=\Sigma_q$. Furthermore, the scheme $C_p\cap\Sigma_q$
has length $2$. 

We prove next that if $C_p\subset \tau_u(\Theta)$, then $\tau_u(\Theta)$ is the unique translate of $\Theta$ containing $C_p\cup \tau_{2u}(\Sigma_p)$. Indeed, the assumed inclusion implies that $u=p-r$, for some $r\in C$, by (\ref{eq-criterion-for-inclusion}). 
Hence, $\tau_u(\Sigma_p)=\Sigma_r\subset\Theta$ and so $\tau_{2u}(\Sigma_p)\subset\tau_u(\Theta)$.
Assume that 
$
C_p\cup \tau_{2u}(\Sigma_p)\subset \tau_u(\Theta)\cap\tau_t(\Theta).
$
Applying $\tau_{-u}$ to both sides we get the inclusion
$
C_r\cup\Sigma_r\subset \Theta\cap \tau_{t-u}(\Theta).
$
Thus, $t-u=r-r=0$, by the previous paragraph. So $t=u$.

\underline{Step 4:} We show next that 
if $C_p\cup\tau_s(\Sigma_p)=\tau_{t_1}(\Theta)\cap\tau_{t_2}(\Theta)$, for some $t_1, t_2\in\Pic^0(C)$, then 
\begin{enumerate}
\item
\label{step-4-item-1}
$t_1=p-q_1$ and $t_2=p-q_2$, for some $q_1, q_2\in C$, 
\item
\label{step-4-item-2}
$s=2p-q_1-q_2$ and $s\neq 0$,
\item
\label{step-4-item-3}
$C_p\cap\tau_s(\Sigma_p)=\left\{p+r, p+t
\right\}$, where $\StructureSheaf{C}(q_1+q_2+r+t)\cong\omega_C$. Furthermore, $C_p$ is tangent to $\tau_s(\Sigma_p)$ at
$p+r$, if $r=t$.
\end{enumerate}
This proves the uniqueness of the pair of translates of $\Theta$ whose complete intersection is $C_p\cup\tau_s(\Sigma_p)$ in Part (\ref{case-a-of-intersection-lemma}).

(\ref{step-4-item-1}) follows from Equation (\ref{eq-criterion-for-inclusion}).

(\ref{step-4-item-2}) 
Applying $\tau_{-t_1}$ we get that 
$C_{q_1}\cup \tau_{s-t_1}(\Sigma_p)=\Theta\cap \tau_{q_1-q_2}(\Theta)$. So $\tau_{s-t_1}(\Sigma_p)=\Sigma_{q_2}$, 
by (\ref{eq-intersection-is-union-of-two-curves}) and $s-t_1=p-q_2$, by (\ref{eq-Sigma-q-is-translate-of-Sigma-p-by-p-q}). This proves the equality in (\ref{step-4-item-2}). If $s=0$, then $p=q_1=q_2$ and so $t_1=t_2$, which contradicts the assumed equality 
$C_p\cup\tau_s(\Sigma_p)=\tau_{t_1}(\Theta)\cap\tau_{t_2}(\Theta)$.

(\ref{step-4-item-3}) We have
$C_{q_1}\cap \tau_{s-t_1}(\Sigma_p)=C_{q_1}\cap \tau_{p-q_2}(\Sigma_p)=C_{q_1}\cap \Sigma_{q_2}=\{q_1+r, q_1+t\}$, 
where $\StructureSheaf{C}(q_1+q_2+r+t)\cong\omega_C$. The first equality follows from (\ref{step-4-item-1})  and (\ref{step-4-item-2}), the second equality from (\ref{eq-Sigma-q-is-translate-of-Sigma-p-by-p-q}), and the third equality from Step 1. If $r=t$ the two curves $C_{q_1}$ and $\Sigma_{q_2}$ are tangent at $q_1+r$, by Step 1. Hence, $C_p$ is tangent to $\tau_s(\Sigma_p)$ at $p+r$.
%and $C_{q_1}\cap \Sigma_{q_2}$ has length $2$. Hence, $C_p\cap\tau_s(\Sigma_p)$ has length $2$. 

%(\ref{step-4-item-4}) 
%We know that $p+r$ and $p+t$ belong to $\tau_s(\Sigma_p)=\tau_s(\iota(C_p))=\iota(\tau_{-s}(C_p)),$ by (\ref{step-4-item-3}).
%Hence, $\tau_s(\iota(p+r))$ belong to $C_p$. Set $L:=\StructureSheaf{C}(s)$. 
%Then $\omega_C(-p-r)\otimes L\cong \StructureSheaf{C}(p+u_1)$, for some $u_1\in C$. 
%Similarly,  $\omega_C(-p-t)\otimes L\cong \StructureSheaf{C}(p+u_2)$, for some $u_2\in C$.
%Thus, 
%\[
%\StructureSheaf{C}(u_1+r)\cong\omega_C\otimes L(-2p)\cong \StructureSheaf{C}(u_2+t).
%\]
%We conclude that $u_1+r=u_2+t$.


%\underline{Step 4(b):} 
%Assume that $C_p\cup \tau_s(\Sigma_p)$ is contained in a unique translate $\tau_u(\Theta)$ of $\Theta$. 
%We have seen in Step 0 that $s=2u$.
%We show that 
%$s=0$ and $u=0$.
%Step 2 implies that $u=p-q$, for some $q\in C$. The inclusion $\tau_{s-u}(\Sigma_p)\subset\Theta$ is equivalent 
%to $\tau_{u-s}(C_p)\subset\Theta$, which implies that $s-u=p-q'$, for some $q'\in C$. 
%So $s=q'-q$ (NO!).  Combined with the equality $s=2u$ we get that $q+q'=2p$. It follows that $p=q=q'$, as $C$ is not hyperelliptic. Thus $s=u=0$.
%**********
% Hide
%**********
\hide{
The following inclusions are equivalent:
\begin{eqnarray*}
C_p\cup \tau_{q'-q}(\Sigma_p) & \subset & \tau_{p-q}(\Theta)=\tau_{p-q}(\iota(\Theta))=\iota(\tau_{q-p}(\Theta))
\\
\Sigma_p\cup \iota(\tau_{q'-q}(\Sigma_p))& \subset &\tau_{q-p}(\Theta)
\\
\Sigma_p\cup \tau_{q-q'}(C_p)& \subset &\tau_{q-p}(\Theta)
\\
\tau_{q'-q}(\Sigma_p)\cup C_p& \subset &\tau_{q'-p}(\Theta)=\tau_{q'+q-2p}(\tau_u(\Theta)).
\end{eqnarray*}
The assumed uniqueness of $u$ implies that $q'+q=2p$. Hence, $p=q'=q$. So $s=0$ and $u=0$. 
The equality $(\tau_u\iota\tau_{-u})(C_p)=\tau_s(\Sigma_p)$ follows.

and $(\tau_u\iota\tau_{-u})(C_p)=\tau_s(\Sigma_p)$. As above we get that $u=p-q$, for some $q\in C$. The inclusion $\tau_{s-u}(\Sigma_p)\subset\Theta$ is equivalent to $\tau_{u-s}(C_p)\subset\Theta$, which implies that $u-s=p-q'$, for some $q'\in C$. Applying $\iota\tau_{-u}$ to the equality 
$(\tau_u\iota\tau_{-u})(C_p)=\tau_s(\Sigma_p)$ and using $\tau_{-u}(C_p)=C_q$ we get
\[
C_q=\iota(\tau_{s-u}(\Sigma_p))=\tau_{u-s}(C_p).
\]
So $p-q'=u-s=q-p$. Thus, $p=q=q'$. So $s=u=0$. This proves the uniqueness of $u$ in 
the ``only if'' direction of part (\ref{lemma-item-length-is-2-iff}) of the lemma if (\ref{case-b-of-intersection-lemma}) holds. 
%**********
% End Hide
%**********
}

%************
% Hide:
%************
\hide{
\underline{Step 5:} Suppose that $C_p\cap\tau_s(\Sigma_p)$ contains a length $2$ subscheme $Z$. So either $Z$ consists of  
two distinct points $\StructureSheaf{C}(p+r)$ and
$\StructureSheaf{C}(p+t)$, or a tangency point $\StructureSheaf{C}(p+r)=\StructureSheaf{C}(p+t)$. 
%If $s=0$ then (\ref{case-b-of-intersection-lemma}) holds for $u=0$. Assume that $s\neq 0$. 
We prove that $C_p\cup\tau_s(\Sigma_p)$ is 
%the intersection of two translates 
contained in a translate
of $\Theta$.
Let $s'\in\Pic^0(C)$ be the point corresponding to $\omega_C^{-1}\otimes  \StructureSheaf{C}(r+t+2p)$.
So $s'=2p-q_1-q_2$, where $q_1+q_2$ is the effective divisor satisfying $\omega_C\cong \StructureSheaf{C}(r+t+q_1+q_2)$.
Then $Z$ in contained also in $C_p\cap \tau_{s'}(\Sigma_p)$, by Step 4(\ref{step-4-item-3}). Furthermore, 
\[
C_p\cup\tau_{s'}(\Sigma_p)\subset \tau_{p-q_1}(\Theta)\cap \tau_{p-q_2}(\Theta).
\]
Indeed, $\tau_{q_1-p}(\tau_{s'}(\Sigma_p))=\tau_{p-q_2}(\Sigma_p)=\Sigma_{q_2}$, by (\ref{eq-Sigma-q-is-translate-of-Sigma-p-by-p-q}), and $\Sigma_{q_2}$  is contained in $\Theta$. Similarly, $\tau_{q_2-p}(\tau_{s'}(\Sigma_p))$   is contained in $\Theta$.
%is the intersection of two translates of $\Theta$, by Step 4. 
We claim that $s=s'$. 
Observe that $\tau_s(\Sigma_p)\cap \tau_{s'}(\Sigma_p)$ contains the subscheme $Z$. Applying $\iota$ we get that 
$\tau_{-s}(C_p)\cap \tau_{-s'}(C_p)$ contains $\iota(Z)$. Hence, $C_p\cap \tau_{s-s'}(C_p)$ contains a length $2$ subscheme. 
We conclude that $s=s'$, by Lemma \ref{lemma-two-translates-of-C}. 
%Now, $s'\neq 0$, by the assumption that $s\neq 0$. 
%We show next that $q_1\neq q_2$, so that $C_p\cup\tau_s(\Sigma_p)$ is the intersection of two translates of $\Theta$. 
%If $q_1=q_2$, then $s=2u$, where $u=p-q_1$. The argument in Step 4(b) shows that $s=0$, which contradicts our assumption.
%************
% End Hide:
%************
}

%\underline{Step 6:} 
%Suppose that $C_p$ and $\tau_s(\Sigma_p)$ are tangent at $\StructureSheaf{C}(p+r).$ The argument in Step 5 still applies, 
%using the case $r=t$ in Step 4, to conclude that $s$ corresponds to the line bundle 
%$\omega_C^{-1}\otimes  \StructureSheaf{C}(2r+2p)$ and to conclude that 
%$C_p\cup \tau_s(\Sigma_p)$ is the intersection of two translates to $\Theta$.

\underline{Step 5:} 
If $C_p\cap \tau_s(\Sigma_p)$ is non-empty, then $C_p\cup \tau_s(\Sigma_p)$ 
is contained in a translate of $\Theta$, by Lemma \ref{lemma-intersection-is-either-empty-or-of-length-2}.
%We have seen that if the subscheme $C_p\cap \tau_s(\Sigma_p)$ has length $\geq 2$, then $C_p\cup \tau_s(\Sigma_p)$ 
%is contained in a translate of $\Theta$. 
It remains to prove that in this case the length of the subscheme $C_p\cap \tau_s(\Sigma_p)$ is $2$. As in the proof of Step 4(3) we can translate so that the subscheme in question is $C_{q_1}\cap \Sigma_{q_2}$ for two points $q_1, q_2\in C$. Now the intersection number of these two curves in $\Theta$ is $2$. This completes the proof of Parts (\ref{lemma-item-length-at-most-2}) and (\ref{lemma-item-length-is-2-iff}).
\end{proof}

\end{document}
